\section{Motivation}

Software engineering is among the most consequential engineering disciplines of the modern era. 
Software systems underpin critical infrastructure across healthcare, finance, transportation, and communication, with failures potentially affecting millions of users, costing billions of dollars, and projecting real-world consequences~\cite{eghbal2016roads, solar-wind, xz-attack}. 
The decisions software engineers make, and the practices they adopt daily (e.g., dependency management strategies, tool adoption choices), accumulate into system-wide outcomes that shape organizational productivity, security posture, and long-term maintainability. 
Given these stakes, one might expect software engineering practice to be grounded in rigorous, evidence-based principles analogous to those in medicine~\cite{sackett1996evidence} or civil engineering~\cite{petroski1985engineer}. 
Unfortunately, this expectation remains largely unfulfilled.

In fact, software engineering practice today is characterized more by folklore, intuition, and contested conventional wisdom than by empirically validated guidance~\cite{DBLP:conf/icse/Devanbu0B16, DBLP:journals/ese/GarousiBO20, DBLP:journals/ese/ShrikanthNFM21, enoiu2026folklore}. 
While practitioners face consequential decisions on a routine basis, they very often find limited reliable evidence to inform their choices. 
To support their decisions, they draw on their own experiences or navigate a landscape of blog posts, podcasts, and conference talks to learn opinions of the ``said'' experts. 
This decision-making process is largely ad hoc, lacks empirical foundations, and may or may not lead to a consensus in the field.
As a result, the field abounds with what I may call \textbf{software engineering myths}, which are \emph{widely held beliefs about best practices derived from intuition, experience, and conventional wisdom, that have never been rigorously tested or, when tested, have produced inconclusive or contradictory results}.

This state of affairs persists despite decades of empirical research in software engineering. 
The field has generated thousands of studies examining key software engineering topics, including developer productivity, code quality, security practices, and the effectiveness of developer tools. 
It has also embraced a vast body of empirical methodological frameworks over the past decades, such as randomized controlled experiments~\cite{DBLP:books/sp/WohlinRHORW24}, discussion of internal/external validity and its trade-offs~\cite{DBLP:conf/icse/SiegmundSA15, DBLP:journals/tosem/RobillardAEGLNNSSS24}, open-science~\cite{DBLP:books/sp/20/0001G0S20}, registered reports~\cite{DBLP:journals/ese/ErnstB23}, qualitative synthesis~\cite{DBLP:conf/esem/CruzesD11}, and more.
This diversity is most apparent in the ACM SIGSOFT Empirical Standards, offering method-specific checklists for 19 major software engineering research methods~\cite{DBLP:journals/corr/abs-2010-03525}.

Despite this substantial body of work, software engineering remains a field with a high degree of folklore, intuition, and conventional wisdom, rather than an evidence-based discipline.
Why is this the case?
There are for sure many unique challenges in the field of empirical software engineering, such as the lack of accurate measurements for key constructs~\cite{DBLP:conf/ease/RalphT18}, the vast ecological diversity of software systems~\cite{DBLP:conf/sigsoft/NagappanZB13}, and the rapid moving nature of software engineering practices in general.
Apart from these challenges, I want to note one important blocker: Many of the software engineering myths we see in practice are \emph{causal in nature}---concerning the effect of an intervention on an outcome of interest---but empirical SE research rarely engages with the challenging question of \emph{whether the presented evidence can actually support a causal claim, and to what extent}.
Instead, researchers often approach software engineering myths with either: (a) descriptive and statistical modeling from software repository data, with a blanket disclaimer that the observed correlation does not imply causation, or (b) qualitative evidence synthesizing and interpreting practitioner opinions on the topic of interest.
Neither would provide sufficient evidence to support or dispel a software engineering myth, especially when it is contested.

For example, consider the following problem: Should software projects allow flexibility in the versions of their dependencies, and to what extent?
A descriptive mining software repository study~\cite{DBLP:conf/msr/0001PSTB19} may conclude that the Java/Maven ecosystem is extensively pinning dependency versions while the JavaScript/npm ecosystem is extensively using semantic versioning, without providing any concrete guidance or discussions of the trade-offs around them.
A qualitative study surveying JavaScript developers~\cite{DBLP:journals/tse/JafariCAST22} may conclude that semantic versioning should be preferred over pinned dependencies as the latter is a ``dependency smell'' likely to cause maintainability issues in the future.
Another qualitative study interviewing security experts~\cite{DBLP:conf/sp/LadisaPMB23} may conclude that dependency pinning should be preferred over semantic versioning, as the former is critical for preventing software supply chain attacks.
In this case, the quantitative repository mining study fails to provide evidence-based guidance, while the two qualitative studies present contradictory opinions from practitioner cohorts with different experiences and perspectives; neither would be sufficient to formulate best practices for this dependency versioning problem.

Similarly, consider the longstanding debate over whether the choice of programming language causally affects defect-proneness. 
A large-scale correlational study~\cite{DBLP:conf/sigsoft/RayPFD14, DBLP:journals/cacm/RayPDF17} analyzed millions of commits across thousands of GitHub projects; the study concluded that programming language has a significant but modest effect on defect-proneness, with functional and statically-typed languages associated with fewer defects.
However, a reproduction study~\cite{DBLP:journals/toplas/BergerHMVV19} revisiting the same dataset with methods that they claim to be better---correcting errors in language classification, accounting for polyglot repositories, and applying more appropriate statistical models---found that the original effect sizes shrank dramatically and many findings failed to replicate. 
The original authors issued a rebuttal~\cite{DBLP:journals/corr/abs-1911-07393} defending their methodology, which prompted a counter-rebuttal~\cite{DBLP:journals/corr/abs-1911-11894} from the reproduction team. 
Putting the controversy around the specific research design aside, the counter-rebuttal~\cite{DBLP:journals/corr/abs-1911-11894} reveals a deeper methodological impasse behind why the field cannot reach consensus on a seemingly straightforward empirical question:
\begin{quote}
\emph{I fundamentally do not believe that the question the FSE authors set out to answer, namely, whether some languages are more defect-prone than others, can be meaningfully answered by scraping GitHub...An analysis of GitHub is necessarily measuring correlation and not establishing causation, as it does not constitute an experiment.}
\end{quote}

The problem extends beyond individual studies and renders empirical software engineering research particularly unproductive on contested problems in which practitioners have strong and conflicting intuitions and beliefs.
When studies only repeat and synthesize practitioner opinions, it may be biased towards a specific perspective and may not be sufficient to shift the posterior beliefs of others.
When studies finds certain correlations from observational data, practitioners will naturally want to read the study as causal without considering the limitations of the study design and the potential confounding factors.
Even worse, when studies reach controversial conclusions (which may be extremely valuable if carefully justified) but fail to rule out important confounding factors or alternative explanations, practitioners may become skeptical of the field altogether. 
The result is a field with prevalent myths despite plentiful empirical research.

Modern causal inference methods offer a principled approach to addressing this problem~\cite{cunningham2021causal, pearl2009causality, hansen2022econometrics}. 
Developed primarily in econometrics and epidemiology over the past several decades, these methods provide frameworks for distinguishing genuine causal effects from spurious correlations in observational data---precisely the kind of data investigated in repository mining studies.
Importantly, the value of causal inference methods lies not in any claim to methodological perfection, but rather in their insistence on \emph{transparency}: Every causal design rests on assumptions that must be explicitly stated, examined, and debated.
When a difference-in-differences study claims a causal effect, it does so under the assumption of parallel trends---an assumption that can be partially tested and openly scrutinized~\cite{marcus2021role}.
When an instrumental variables analysis identifies an effect, it does so under exclusion restrictions that critics can directly challenge~\cite{deaton2010instruments}.
This methodological transparency addresses the core problem illustrated by the programming language vs. defect-proneness example: 
When studies produce contested results, progress requires not endless reanalysis but a shared framework for evaluating whether the findings and claims are plausible given the specific research context and, if not, how to develop a better design.

In this thesis, I argue that \emph{causal credibility assessments are feasible in empirical software engineering through explicit articulations of causal theory and assumptions}.
This pragmatic stance recognizes that causal inference in observational settings is inherently imperfect: No study eliminates all confounding, no natural experiment is as clean as a randomized controlled trial, and no finding generalizes to all contexts without qualification.
Yet, these limitations do not render causal inquiry futile.
Rather, they demand intellectual honesty about what we can and cannot conclude from the available data sources, and careful calibration of confidence to convey the strength of the evidence to the audience in each study.

Explicit causal credibility assessments are particularly important for research challenging well-known software engineering myths.
From a Bayesian perspective, if a study yields results contrary to practitioners' intuitions, it requires more compelling evidence to shift posterior beliefs than a study that merely confirms conventional wisdom.\footnote{This reflects the Bayesian epistemological principle often summarized as ``extraordinary claims require extraordinary evidence.'' When a claim has a low prior probability (i.e., it contradicts established beliefs or practitioner intuition), stronger evidence is needed to shift the posterior belief toward acceptance. Formally, if $P(H)$ is low, then $P(E|H)/P(E|\neg H)$---the likelihood ratio---must be correspondingly high for $P(H|E)$ to reach credible levels. For a rigorous treatment of Bayesian reasoning in scientific inference, see Howson and Urbach~\cite{howson2006scientific}.}
The value of explicit causal credibility assessments lies not in the capacity to provide causal evidence of epistemological certainty, but in the transparent articulation of theory and assumptions, the honest acknowledgment and mitigation of limitations, and the careful accumulation of consistent findings across complementary designs.
The three empirical studies in this thesis (Chapter~\ref{chp:pinning}, \ref{chp:fake-stars}, \ref{chp:cursor}) exemplify this approach: Each challenges a widely held belief and explicitly assesses the credibility of its causal claims.

\section{Thesis Statement}

\begin{framed}
%Software engineering practice is characterized by \emph{myths}---widely held beliefs about best practices derived from intuition, experience, and conventional wisdom that have never been rigorously tested or, when tested, have produced inconclusive results.
In this thesis, I argue that \emph{causal credibility assessments is feasible in empirical software engineering through an explicit articulation of causal theory and assumptions}.
\end{framed}

There are multiple established approaches in other disciplines for causal credibility assessments---from Bradford Hill's viewpoints for evaluating epidemiological evidence~\cite{hill1965environment}, to the design-based credibility revolution in empirical economics that emphasizes transparent identification strategies~\cite{angrist2010credibility}, to sensitivity analysis frameworks that quantify robustness to unmeasured confounding~\cite{rosenbaum2002observational}.
Drawing on these traditions, this thesis demonstrates the feasibility of causal credibility assessments in software engineering through one particular operationalization: a four-step process that I apply consistently across the empirical studies in this thesis:
\begin{enumerate}[leftmargin=20pt]
    \item Derive an explicit causal theory from existing software engineering domain knowledge---for instance, as a directed acyclic graph (DAG), which is especially valuable for regression-based studies where variable selection determines the credibility of causal claims,
    \item Define the causal estimate of interest through the potential outcome framework and specify the assumption under which the estimate would carry a causal interpretation,
    \item Honestly acknowledge and mitigate limitations and known caveats in data collection, metric design, measurement infrastructure, and model specifications, 
    \item Thoughtfully navigate alternative explanations (i.e., alternative causal structures or causal DAGs) that may also be supported with the current results. 
\end{enumerate}
This four-step process is not the only way to conduct causal credibility assessments, nor does this thesis claim it is superior to alternative approaches.
Other operationalizations---emphasizing different combinations of formal sensitivity analysis~\cite{rosenbaum2002observational, vanderweele2017sensitivity}, structured reporting guidelines, or pre-registered analysis plans---may be equally valid depending on the research context.
The contribution of this thesis is to demonstrate that the proposed form of explicit causal credibility assessment is both feasible and productive in the software engineering research context, where such assessments have been largely absent.
The three empirical studies (Chapters~\ref{chp:pinning}, \ref{chp:fake-stars}, and~\ref{chp:cursor}) illustrate this feasibility by applying these four steps to three distinct research problems, each challenging a widely held software engineering myths.

\section{Organization and Contributions}

In this section, I summarize the organization and main contributions of this thesis.

\paragraph{A Tutorial on Causal Inference for Empirical Software Engineering (Chapter~\ref{chp:background}).}
The thesis begins with a self-contained tutorial on causal inference for empirical software engineering researchers.
The chapter opens by tracing the origins and scope of the modern causal inference toolkit, surveying existing methodological guides in software engineering, and documenting the toolkit's limited adoption in the field to date.
It then develops a primer that articulates the fundamental problem of causal inference---the impossibility of observing both potential outcomes for a single unit---and explains why descriptive and correlational evidence, including multivariate regression, cannot on its own establish causal claims.
The primer introduces the three conceptual frameworks that have shaped modern causal reasoning---the potential outcomes framework, graphical causal models, and the design-based approach---and synthesizes them into a pragmatic stance on causal credibility, operationalizing the four-step process introduced in the thesis statement and articulating the internal--external validity trade-off that any identification strategy must navigate.
To demonstrate the toolkit in practice, the chapter then walks through a worked example reanalyzing the impact of AI coding tools on open-source projects, contrasting cross-sectional and longitudinal designs and showing how each design's identifying assumptions shape the credibility of its conclusions.
The chapter closes by drawing lessons from the worked example, discussing implications for empirical software engineering, articulating the structural advantages that software engineering offers for causal inference relative to other disciplines, and identifying limitations and open frontiers.

\paragraph{Three Cases of Empirical Studies.}
Then, the thesis dives deep into three empirical studies I published in the past that challenged well-known software engineering myths through a pragmatic causal inference with explicit causal credibility assessments---each study deriving an explicit causal theory from domain knowledge, defining the causal estimate of interest and specifying the identifying assumption, honestly acknowledging and mitigating limitations and known caveats, and thoughtfully navigating alternative explanations~\cite{DBLP:journals/corr/abs-2502-06662, DBLP:journals/corr/abs-2412-13459, DBLP:journals/corr/abs-2511-04427}.

\begin{itemize}[leftmargin=15pt]
    \item \textbf{The Impact of Dependency Pinning on Supply Chain Security (Chapter~\ref{chp:pinning})}: Recent high-profile incidents in open-source software have greatly increased practitioner attention on software supply chain attacks. 
    To guard against potential malicious package updates, security practitioners advocate \emph{pinning} dependency to specific versions rather than \emph{floating} in version ranges.
    However, it remains controversial whether pinning provides a meaningful security benefit that outweighs the cost of maintaining outdated, possibly vulnerable dependencies.
    In this chapter, I quantify, through counterfactual analysis and simulations, the security and maintenance impact of version constraints in the npm ecosystem. 
    By simulating dependency resolutions over historical time points, I find that pinning direct dependencies not only (as expected) increases the cost of maintaining vulnerable and outdated dependencies, but also (surprisingly) even increases the risk of exposure to malicious package updates in larger dependency graphs due to the specifics of npm's dependency resolution mechanism. 
    Finally, I explore collective pinning strategies to secure the ecosystem against supply chain attacks and suggest specific changes to npm to enable these interventions.
    These findings provide guidance to practitioners and tool designers on how to manage their supply chains more securely.
    \item  \textbf{The Impact of Fake GitHub Stars on Repository Promotion (Chapter~\ref{chp:fake-stars})}: GitHub, the de facto platform for open-source software development, provides a set of social-media-like features to signal high-quality repositories.
    Among them, the star count is the most widely used popularity signal, but it is also at risk of being artificially inflated (i.e., \emph{faked}), decreasing its value as a decision-making signal and posing a security risk to all GitHub users.
    In this chapter, I present a systematic, global, and longitudinal measurement study of fake stars in GitHub.
    To this end, I built \SystemName, a scalable tool able to detect anomalous starring behaviors across all GitHub metadata between 2019 and 2024.
    Analyzing the data collected using \SystemName, I find that: (1) fake-star-related activities have rapidly surged in 2024; (2) the accounts and repositories in fake star campaigns have highly trivial activity patterns; (3) the majority of fake stars are used to promote short-lived phishing malware repositories; the remaining ones are mostly used to promote AI/LLM, blockchain, tool/application, and tutorial/demo repositories; (4) while repositories may have acquired fake stars for growth hacking, fake stars only have a promotion effect in the short term (i.e., less than two months) and become a liability in the long term.
    This study has implications for platform moderators, open-source practitioners, and supply chain security researchers.
    \item \textbf{The Impact of Cursor AI on Open-Source Projects (Chapter~\ref{chp:cursor})}:   Large language models (LLMs) have demonstrated the promise to revolutionize the field of software engineering.
    Among other things, LLM agents are rapidly gaining momentum in software development, with practitioners reporting a multifold increase in productivity after adoption.
    Yet, empirical evidence is lacking around these claims.
    In this chapter, I estimate the \emph{causal} effect of adopting a widely popular LLM agent assistant, namely Cursor, on \emph{development velocity} and \emph{software quality}.
    The estimation is enabled by a state-of-the-art difference-in-differences design comparing Cursor-adopting GitHub projects with a matched control group of similar GitHub projects that do not use Cursor. 
    I find that Cursor adoption leads to a significant, large, but transient increase in project-level development velocity, along with a substantial and persistent increase in static analysis warnings and code complexity.
    Further panel generalized method of moments estimation reveals that the increase in static analysis warnings and code complexity acts as a major factor causing long-term velocity slowdown.
    This study carries implications for software engineering practitioners, LLM agent assistant designers, and researchers. 
\end{itemize}

\paragraph{Summary of Contributions.} This thesis makes the following contributions:

\begin{itemize}[leftmargin=15pt]
    \item \textbf{Conceptual Contribution.} Many software engineering best practices are widely held but lack rigorous empirical validation---a phenomenon I call \emph{software engineering myths}.
    I argue that their persistence reflects a structural feature of the field: the reliance on a blanket disclaimer ``correlation does not imply causation'' and the absence of an accessible causal inference infrastructure for conveying the extent to which the evidence can support a causal claim in each study.
    Against this diagnosis, I propose \emph{explicit causal credibility assessment} as a pragmatic methodological standard for empirical software engineering, prioritizing the transparent disclosure of assumptions over epistemological perfection.
    The standard is operationalized as a four-step process: (1) derivation of an explicit causal theory from existing software engineering domain knowledge (e.g., as DAGs, which is especially valuable for regression-based studies where variable selection determines the credibility of causal claims); (2) definition of the causal estimate of interest through the potential outcomes framework and specification of the assumption under which the estimate carries a causal interpretation; (3) honest acknowledgment and mitigation of limitations and known caveats in data collection, metric design, measurement infrastructure, and model specifications; and (4) thoughtful navigation of alternative explanations (i.e., alternative causal structures or causal DAGs) that may also be supported by the current results.
    Together, the diagnosis and the standard provide a path for incremental accumulation of trustworthy evidence in software engineering, a lens for critically evaluating existing empirical work, and a principled guideline for future studies.

    \item \textbf{Methodological Contributions.} Existing causal inference methods rarely apply off-the-shelf to software engineering data.
    This thesis contributes three reusable design adaptations that bridge SE-specific data structures and standard causal designs, each carrying its own credibility assessment in the sense above:
    \begin{enumerate}[leftmargin=15pt, label=(\alph*)]
        \item \emph{Counterfactual simulation of dependency resolution at ecosystem scale}, exploiting an undocumented time-travel feature of the npm resolver to construct same-project treatment/control panels across historical time points, paired with fixed-effects panel regression interacted with dependency-graph size to recover effect heterogeneity that prior descriptive studies could not detect (Chapter~\ref{chp:pinning}).
        \item \emph{Anomaly-detection signatures for coordinated platform manipulation on GitHub}, adapting lockstep- and low-activity-pattern detection from social-network fraud detection to the full GitHub event archive, together with a deletion-rate-elevation protocol for validating such detectors when ground truth is partial (Chapter~\ref{chp:fake-stars}).
        \item \emph{A quasi-experimental design template for studying developer-tool adoption in open-source software}, combining configuration-file-based treatment identification from git history, propensity score matching on pre-adoption \emph{dynamics} rather than static snapshots, the \citet{borusyak2024revisiting} staggered-adoption imputation estimator, and panel GMM for recovering bidirectional dynamics between mediators and outcomes (Chapter~\ref{chp:cursor}).
    \end{enumerate}

    \item \textbf{Empirical Contributions.} This thesis produces novel causal evidence overturning three widely held software engineering myths:
    \begin{enumerate}[leftmargin=15pt, label=(\alph*)]
        \item \emph{Dependency pinning is not a uniformly safer default for the npm ecosystem.}
        Counterfactual evidence from a 20{,}000-project, 12-month panel shows that pinning direct dependencies increases exposure to malicious updates in larger dependency graphs (becoming actively counterproductive at $\geq$498 dependencies, affecting 26\% of GitHub repositories in the dataset), while simultaneously increasing exposure to known vulnerabilities, outdated dependencies, and dependency-graph bloat.
        Simulation further shows that coordinated pinning of carefully selected upstream packages reduces ecosystem-wide malicious-update risk by up to 75\%---far exceeding what individual pinning can achieve---motivating concrete proposed changes to npm to enable such coordination (Chapter~\ref{chp:pinning}).
        \item \emph{Fake GitHub stars are a long-term liability, not a sustainable promotional asset.}
        The first global, longitudinal characterization of fake-star campaigns (6.0M suspected fake stars across 18{,}617 repositories in 2019--2024, with peak-month prevalence of 16.66\% among popular repositories in 2024) reveals that the dominant use case is short-lived phishing-malware promotion, with the remainder concentrated in AI/LLM, blockchain, tooling, and tutorial repositories.
        Panel autoregression estimates show fake stars have a short-term promotional effect roughly five times smaller than genuine stars and a \emph{negative} cumulative effect after two months (Chapter~\ref{chp:fake-stars}).
        \item \emph{LLM coding agents trade transient velocity gains for persistent quality degradation in real open-source projects.}
        The first large-scale causal estimate of Cursor adoption on GitHub finds a 3--5$\times$ increase in lines added in the first month post-adoption that dissipates within two months, alongside a persistent +30\% increase in static-analysis warnings and +41\% increase in cognitive complexity that does not recover over the observation window.
        Panel GMM further identifies a self-reinforcing feedback loop in which accumulated technical debt causally reduces subsequent velocity, providing a mechanism for the dissipation of the initial productivity surge (Chapter~\ref{chp:cursor}).
    \end{enumerate}
\end{itemize}